\section{Price Scope and Adaptation Checks}
\label{app:scope_adaptation_v18}

The main text treats price evidence as corroboration rather than
identification. This appendix records the two checks that keep
that interpretation disciplined: the overlap sign reversal and
the adversarial-adaptation simulation.

\subsection{Overlap and Segment Scope}
\label{app:overlap_scope_v18}

The broad price association is positive, with the headline range
$\valHeadlineRange$. Under overlap-cell ATT reweighting, the
aggregate coefficient reverses to $\valScopeOverlapCoef$; under
propensity-score trimming it becomes $\valScopePSCoef$. This
does not refute the screening interpretation. It shows that the
broad estimate and the overlap estimate average different parts
of the procurement distribution.

The segment decomposition is the useful object. The lowest three
tender-value quartiles have negative coefficients
($\valSRQoneBroad$, $\valSRQtwoBroad$, and
$\valSRQthreeBroad$ in the broad specification). The highest
tender-value quartile is positive in both the broad and overlap
designs: $\valSRQfourBroad$ and $\valSRQfourATT$
($p=\valSRQfourATTp$). Trimming heavily weighted overlap cells
does not eliminate the negative aggregate estimate; it moves from
$\valSRTrimNone$ with no trim to $\valSRTrimTen$ after dropping
the top decile of cell weights and $\valSRTrimFifty$ after
dropping the top half.

The legal consequence is narrow. The price results show that
flagged environments are economically non-neutral, especially in
the high-value segment where cartel rents are most plausible.
They do not supply a causal price effect, a markup, or a damages
calculation.

\subsection{Adversarial Adaptation}
\label{app:adaptation_v18}

The screen is not invariant to strategic response. The
simulation holds the empirical environment fixed and asks how the
participation score behaves under simple evasive changes. Rotation
of cover bidders leaves AUC essentially unchanged
($\valAUCAdaptRotation$ versus baseline $\valAUCBase$).
Occasional wins also leave AUC nearly unchanged
($\valAUCAdaptOccWins$), because the strict zero-win rule removes
those firms from the flagged set. CNPJ splitting lowers AUC to
$\valAUCAdaptCNPJsplit$. Threshold-aware capping is more damaging,
lowering AUC to $\valAUCAdaptThreshCap$. Combining capping with
rotation lowers AUC to $\valAUCAdaptCombined$.

\begin{table}[htbp]
\centering
\caption{Adversarial-Adaptation Summary}
\label{tab:adaptation_summary_v18}
\footnotesize
\begin{adjustbox}{max width=\textwidth,center}
\begin{tabular}{p{4.4cm}cc}
\toprule
Scenario & AUC & Change vs baseline \\
\midrule
Baseline & $\valAUCBase$ & --- \\
Rotation & $\valAUCAdaptRotation$ & $\valAdaptRotationDeltaPP$ pp \\
Occasional wins & $\valAUCAdaptOccWins$ & $\valAdaptOccWinsDeltaPP$ pp \\
CNPJ splitting & $\valAUCAdaptCNPJsplit$ & $\valAdaptCNPJDeltaPP$ pp \\
Threshold capping & $\valAUCAdaptThreshCap$ & $\valAdaptCapDeltaPP$ pp \\
Capping plus rotation & $\valAUCAdaptCombined$ & $\valAdaptCombinedDeltaPP$ pp \\
\bottomrule
\end{tabular}
\end{adjustbox}
\end{table}

The practical lesson is not to hide the cutoff. It is to treat
the cutoff as operational rather than structural. Agencies should
refresh thresholds, monitor bunching below the rule, and pair the
award-layer screen with bid-layer forensics when bid microdata
can be recovered. Strategic response weakens a static binary rule
more than it weakens the two-stage architecture.
